% Chapter 1 - Getting Started.

\guidechapter{GETTING STARTED}{%
  \item What the SDK requires
  \item Enabling the command interface
  \item Choosing and checking a binding
  \item Recovering from common failures
}

\section{What The SDK Requires}

The SDK runs on a C64 with Ultimate hardware whose command interface is mapped
at \cmd{\$DF1B--\$DF1F}. It is supported on the 1541~Ultimate-II and later
products.

There is no blanket minimum firmware version. Targets were added over time, so
programs call \cmd{ultimate\_detect()} and test the capability they need. The
palette is newer than firmware 3.15; HTTP appears in 3.15.

The control target predates the palette commands, so its presence is not a
palette capability test. Call a palette function and handle
\cmd{ULTIMATE\_ERR\_NOT\_SUPPORTED}; the SDK returns that result when older
firmware answers that it does not know the command.

\evidence{\cmd{docs/compatibility.md}; \cmd{docs/uci.md};
\cmd{include/ultimate.h}}

\section{Enable The Command Interface}

Open the Ultimate configuration menu and set \bas{command interface} to
\bas{enabled}. The exact key used to open that menu depends on the product and
its configuration, so the SDK does not prescribe one.

With the interface disabled, the five registers are not mapped and
\cmd{ultimate\_init()} returns \cmd{ULTIMATE\_ERR\_NO\_DEVICE}.

REU transfers and Ultimate~64 turbo control use their own registers, not the
command interface. Each still depends on its corresponding machine setting.

\evidence{\cmd{docs/compatibility.md};
\cmd{tests/hardware/hwtest.py}; \cmd{tools/test\_registers.py}}

\section{Choose A Binding}

\begin{reftable}{lp{0.63\linewidth}}
BASIC V2 & \cmd{make wedge}; load \cmd{src/basic/uci.prg}, or mount
  \cmd{src/basic/uci.crt} \\
cc65 C & \cmd{make lib}; include \cmd{ultimate.h} and link
  \cmd{bindings/cc65/build/ultimate.lib} \\
ca65 assembly & link the same library and include
  \cmd{bindings/asm/ultimate.inc} \\
other toolchains & \cmd{make blob}; load the standalone binary described in
  Chapter~5 \\
\end{reftable}

\evidence{\cmd{Makefile}; \cmd{src/basic/Makefile};
\cmd{bindings/cc65/Makefile}; \cmd{bindings/blob/Makefile}}

\Needspace{18\baselineskip}
\subsection*{BASIC check}

Load and run the wedge, list the current directory, then print the command
result:

\begin{screen}
LOAD"UCI",8\\
RUN\\
UDIR\\
PRINT UERR\\
0
\end{screen}

This demonstrates wedge installation, a DOS command and the shared result
model. A printed \bas{0} is \cmd{ULTIMATE\_OK}; \bas{1} is
\cmd{ULTIMATE\_ERR\_NO\_DEVICE}.

\evidence{\cmd{src/basic/loader.s}; \cmd{src/basic/dispatch.s};
\cmd{include/uci\_protocol.h}}

\subsection*{C check}

This complete program demonstrates initialization and error reporting:

\begin{listingbox}
\#include <stdio.h>\\
\#include <ultimate.h>\\
\\
int main(void)\\
\{\\
\hspace*{1em}uint8\_t err = ultimate\_init();\\
\\
\hspace*{1em}if (err != ULTIMATE\_OK)\\
\hspace*{2em}printf("ultimate: \%s\textbackslash n", ultimate\_strerror(err));\\
\hspace*{1em}return err != ULTIMATE\_OK;\\
\}
\end{listingbox}

\begin{outputbox}
No text on success. On failure, for example:\\
ULTIMATE: NO ULTIMATE FOUND
\end{outputbox}

Build it from the repository root:

\begin{listingbox}
cl65 -t c64 -I include -o hello.prg hello.c \textbackslash\\
\hspace*{2em}bindings/cc65/build/ultimate.lib
\end{listingbox}

\begin{outputbox}
hello.prg is created with no compiler errors.
\end{outputbox}

\evidence{\cmd{include/ultimate.h}; \cmd{examples/cc65/Makefile}}

\Needspace{14\baselineskip}
\subsection*{Assembly check}

This demonstrates the assembly return convention: the result code is in
\cmd{A}, ready to compare and branch on.

\begin{listingbox}
\hspace*{2em}jsr ultimate\_init\\
\hspace*{2em}cmp \#ULTIMATE\_OK\\
\hspace*{2em}bne failed
\end{listingbox}

\begin{outputbox}
A = ULTIMATE_OK; BNE is not taken.
\end{outputbox}

The \cmd{cmp} instruction leaves the processor flags ready for \cmd{bne}.

\evidence{\cmd{docs/asm-abi.md}; \cmd{examples/asm/identify.s}}

\section{Recover From Common Failures}

\subsection*{No device}

\cmd{ULTIMATE\_ERR\_NO\_DEVICE} means the signature is absent. Check the
Command Interface setting before assuming the hardware is missing.

\Needspace{12\baselineskip}
\subsection*{An abandoned exchange}

A program can stop while the firmware is holding command or reply state.
\cmd{uci\_abort()} releases it; \cmd{ultimate\_init()} performs that recovery
during startup. From BASIC, a bare \bas{uci} calls the same abort path:

\begin{screen}
UCI:PRINT UERR\\
0\\
READY.
\end{screen}

This example demonstrates recovery only; it does not send a command.

\subsection*{A missing feature}

Use \cmd{ultimate\_detect()} for command-interface targets,
\cmd{ultimate\_reu\_available()} for the REU, and
\cmd{ultimate\_turbo\_available()} for turbo. Do not infer capabilities from a
model name.

\subsection*{A short buffer}

\cmd{ULTIMATE\_ERR\_TRUNCATED} means an exchange completed but not all reply
bytes fitted. Output lengths report the bytes retained, not the bytes discarded.

\evidence{\cmd{include/uci.h}; \cmd{include/ultimate.h};
\cmd{src/basic/dispatch.s}; \cmd{tests/emulator/sdk.suite}}
